% GENERATED FILE. Edit canonical lessons, then run npm run generate:books.
% canonical-source-hash: fnv1a64:2cf7d7aa18e826d3
% canonical-lessons: BN-C231-jhamela, BN-C231-somadhan, BN-C231-ata2, BN-C231-moyda, BN-C231-suji, BN-R231-first-pass-words-from-chapters-227-228, BN-R231-second-pass-words-from-chapters-229-231

\chapter{Hassle, Solution, Wholemeal flour, Fine flour, Semolina}
\label{ch:bn-hassle-solution-wholemeal-flour-fine-flour-semolina}
\hlchaptermodality{231}

\begin{chapteropening}
\textbf{By the end of this chapter:} \emph{I can say a hassle, a solution, wholemeal flour, fine flour and semolina.}

Complete the last lesson of chapter 231: I can say a hassle, a solution, wholemeal flour, fine flour and semolina.
\end{chapteropening}

\section[\texorpdfstring{\bn{ঝামেলা} (jhāmelā)}{ঝামেলা (jhāmelā)}]{\bn{ঝামেলা} (jhāmelā) — a hassle}

\label{lesson:BN-C231-jhamela}



\begin{quote}
Before the new one: say the Bengali for society, then the Bengali for religion.
\end{quote}

\subsection*{You'll want to know: \bn{ঝামেলা}}
\textbf{\bn{ঝামেলা}} — \emph{jhāmelā} — \textquotedblleft{}a hassle\textquotedblright{}.

\begin{grammarlens}[title={What you've built}]
One more word for this chapter.
\end{grammarlens}

\subsection*{Your turn}
\begin{itemize}
\raggedright
  \item \emph{Say it:} \emph{jhāmelā}
  \item \emph{Say it:} \emph{jhāmelā}, once more
  \item \emph{Say it:} say \emph{jhāmelā}
  \item \emph{Recall:} say the Bengali for society, then the Bengali for religion, then say \emph{jhāmelā} again
\end{itemize}

\begin{tcolorbox}[breakable,colback=teal!4,colframe=teal!35!black,title={Before you move on}]
What does \bn{ঝামেলা} mean? (\textquotedblleft{}A hassle\textquotedblright{}.) Say it once more.
\end{tcolorbox}
\section[\texorpdfstring{\bn{সমাধান} (sômādhān)}{সমাধান (sômādhān)}]{\bn{সমাধান} (sômādhān) — a solution}

\label{lesson:BN-C231-somadhan}



\begin{quote}
Before the new one: say the Bengali for religion, then the Bengali for a hassle.
\end{quote}

\subsection*{You'll want to know: \bn{সমাধান}}
\textbf{\bn{সমাধান}} — \emph{sômādhān} — \textquotedblleft{}a solution\textquotedblright{}.

\begin{grammarlens}[title={What you've built}]
One more word for this chapter.
\end{grammarlens}

\subsection*{Your turn}
\begin{itemize}
\raggedright
  \item \emph{Say it:} \emph{sômādhān}
  \item \emph{Say it:} \emph{sômādhān}, once more
  \item \emph{Say it:} say \emph{sômādhān}
  \item \emph{Recall:} say the Bengali for religion, then the Bengali for a hassle, then say \emph{sômādhān} again
\end{itemize}

\begin{tcolorbox}[breakable,colback=teal!4,colframe=teal!35!black,title={Before you move on}]
What does \bn{সমাধান} mean? (\textquotedblleft{}A solution\textquotedblright{}.) Say it once more.
\end{tcolorbox}
\section[\texorpdfstring{\bn{আটা} (āṭā)}{আটা (āṭā)}]{\bn{আটা} (āṭā) — wholemeal flour}

\label{lesson:BN-C231-ata2}



\begin{quote}
Before the new one: say the Bengali for a hassle, then the Bengali for a solution.
\end{quote}

\subsection*{You'll want to know: \bn{আটা}}
\textbf{\bn{আটা}} — \emph{āṭā} — \textquotedblleft{}wholemeal flour\textquotedblright{}.

\begin{grammarlens}[title={What you've built}]
One more word for this chapter.
\end{grammarlens}

\subsection*{Your turn}
\begin{itemize}
\raggedright
  \item \emph{Say it:} \emph{āṭā}
  \item \emph{Say it:} \emph{āṭā}, once more
  \item \emph{Say it:} say \emph{āṭā}
  \item \emph{Recall:} say the Bengali for a hassle, then the Bengali for a solution, then say \emph{āṭā} again
\end{itemize}

\begin{tcolorbox}[breakable,colback=teal!4,colframe=teal!35!black,title={Before you move on}]
What does \bn{আটা} mean? (\textquotedblleft{}Wholemeal flour\textquotedblright{}.) Say it once more.
\end{tcolorbox}
\section[\texorpdfstring{\bn{ময়দা} (môydā)}{ময়দা (môydā)}]{\bn{ময়দা} (môydā) — fine flour}

\label{lesson:BN-C231-moyda}



\begin{quote}
Before the new one: say the Bengali for a solution, then the Bengali for wholemeal flour.
\end{quote}

\subsection*{You'll want to know: \bn{ময়দা}}
\textbf{\bn{ময়দা}} — \emph{môydā} — \textquotedblleft{}fine flour\textquotedblright{}.

\begin{grammarlens}[title={What you've built}]
One more word for this chapter.
\end{grammarlens}

\subsection*{Your turn}
\begin{itemize}
\raggedright
  \item \emph{Say it:} \emph{môydā}
  \item \emph{Say it:} \emph{môydā}, once more
  \item \emph{Say it:} say \emph{môydā}
  \item \emph{Recall:} say the Bengali for a solution, then the Bengali for wholemeal flour, then say \emph{môydā} again
\end{itemize}

\begin{tcolorbox}[breakable,colback=teal!4,colframe=teal!35!black,title={Before you move on}]
What does \bn{ময়দা} mean? (\textquotedblleft{}Fine flour\textquotedblright{}.) Say it once more.
\end{tcolorbox}
\section[\texorpdfstring{\bn{সুজি} (suji)}{সুজি (suji)}]{\bn{সুজি} (suji) — semolina}

\label{lesson:BN-C231-suji}



\begin{quote}
Before the new one: say the Bengali for wholemeal flour, then the Bengali for fine flour.
\end{quote}

\subsection*{You'll want to know: \bn{সুজি}}
\textbf{\bn{সুজি}} — \emph{suji} — \textquotedblleft{}semolina\textquotedblright{}.

\begin{grammarlens}[title={What you've built}]
One more word for this chapter.
\end{grammarlens}

\subsection*{Your turn}
\begin{itemize}
\raggedright
  \item \emph{Say it:} \emph{suji}
  \item \emph{Say it:} \emph{suji}, once more
  \item \emph{Say it:} say \emph{suji}
  \item \emph{Recall:} say the Bengali for wholemeal flour, then the Bengali for fine flour, then say \emph{suji} again
\end{itemize}

\begin{tcolorbox}[breakable,colback=teal!4,colframe=teal!35!black,title={Before you move on}]
What does \bn{সুজি} mean? (\textquotedblleft{}Semolina\textquotedblright{}.) Say it once more.
\end{tcolorbox}
\section[(dialogue)]{Words from chapters 227-228 — a first pass}

\label{lesson:BN-R231-first-pass-words-from-chapters-227-228}



\begin{quote}
Say the words for fine flour and semolina.
\end{quote}

\subsection*{Your turn}
Say these aloud:

\begin{itemize}
\raggedright
  \item a surname, an age, a number, advice, the truth
  \item failure, life, death, luck, a decision
\end{itemize}

\begin{tcolorbox}[breakable,colback=teal!4,colframe=teal!35!black,title={Before you move on}]
A surname? (\textbf{\bn{পদবি}}, \emph{pôdôbi}.) An age? (\textbf{\bn{বয়স}}, \emph{bôyôs}.) A number? (\textbf{\bn{নম্বর}}, \emph{nômbôr}.) Advice? (\textbf{\bn{পরামর্শ}}, \emph{pôrāmôrshô}.) The truth? (\textbf{\bn{সত্য}}, \emph{sôtyô}.) Failure? (\textbf{\bn{ব্যর্থতা}}, \emph{bærthôtā}.) Life? (\textbf{\bn{জীবন}}, \emph{jībôn}.) Death? (\textbf{\bn{মৃত্যু}}, \emph{mrityu}.) Luck? (\textbf{\bn{ভাগ্য}}, \emph{bhāgyô}.) A decision? (\textbf{\bn{সিদ্ধান্ত}}, \emph{siddhāntô}.)
\end{tcolorbox}
\section[(dialogue)]{Words from chapters 229-231 — a second pass}

\label{lesson:BN-R231-second-pass-words-from-chapters-229-231}



\begin{quote}
Say the words for a torch and a chain.
\end{quote}

\subsection*{Your turn}
Say these aloud:

\begin{itemize}
\raggedright
  \item a torch, a chain, a saw, a ribbon, the weather
  \item a leader, a minister, a prince, society, religion
  \item a hassle, a solution, wholemeal flour, fine flour, semolina
\end{itemize}

\begin{tcolorbox}[breakable,colback=teal!4,colframe=teal!35!black,title={Before you move on}]
A torch? (\textbf{\bn{টর্চ}}, \emph{ṭôrch}.) A chain? (\textbf{\bn{শিকল}}, \emph{shikôl}.) A saw? (\textbf{\bn{করাত}}, \emph{kôrāt}.) A ribbon? (\textbf{\bn{ফিতে}}, \emph{phite}.) The weather? (\textbf{\bn{আবহাওয়া}}, \emph{ābôhāoyā}.) A leader? (\textbf{\bn{নেতা}}, \emph{netā}.) A minister? (\textbf{\bn{মন্ত্রী}}, \emph{môntrī}.) A prince? (\textbf{\bn{রাজকুমার}}, \emph{rājkumār}.) Society? (\textbf{\bn{সমাজ}}, \emph{sômāj}.) Religion? (\textbf{\bn{ধর্ম}}, \emph{dhôrmô}.) A hassle? (\textbf{\bn{ঝামেলা}}, \emph{jhāmelā}.) A solution? (\textbf{\bn{সমাধান}}, \emph{sômādhān}.) Wholemeal flour? (\textbf{\bn{আটা}}, \emph{āṭā}.) Fine flour? (\textbf{\bn{ময়দা}}, \emph{môydā}.) Semolina? (\textbf{\bn{সুজি}}, \emph{suji}.)
\end{tcolorbox}
